% Thesis abstract, kept in its own file for the submission system, which
% ingests the abstract separately from the manuscript (Guide §10.3).
% Limit: 3,500 characters. The build does not input this file.

Score-based diffusion models generate data by learning, at every noise
level, the gradient of the log-density of progressively corrupted data ---
the score --- and integrating a reverse-time dynamics driven by it. For the
sequential data these models increasingly target, the score is almost
always treated as a black box to be approximated by a neural network. This
thesis studies a setting in which nothing about it needs to be guessed:
sequences generated by a first-order Markov chain and corrupted by the
variance-preserving diffusion channel, where the joint score can be
characterised exactly, the cost of approximating it can be measured against
a validated reference, and the value of knowing the generative structure
can be quantified against networks that are not told it.

Three groups of results are established. First, structure: for the Gaussian
AR(1) chain the joint score is linear, $S(x,t) = -Q_t x$, and the precision
$Q_t$ has a closed-form lifecycle --- exactly tridiagonal at $t=0$, filling
one diagonal per order in $t$ with every off-band coupling provably
non-zero at all $t>0$, and returning to the identity at rate $e^{-2t}$; a
single bulk quantity $q(\alpha,t)$ sets the posterior correlation length,
and the influence of evidence at distance $d$ decays as $q^d$. Belief
propagation on the noised chain is shown to be, update by update, the
Kalman filter and Rauch--Tung--Striebel smoother, so one forward and one
backward sweep deliver every marginal exactly.

Second, approximation: when the innovations are Laplace rather than
Gaussian, the messages no longer close in finite dimension. The thesis
proves that moment-projected (Gaussian) message passing returns precisely
the linear minimum-mean-square-error estimator of the covariance-matched
Gaussian model, turning the closure ``approximation error'' into a
well-defined price of discarding information beyond second moments;
against a resolution-validated grid reference this price is measured at
median relative score error $0.27$ at $t=0.05$, falling below $10^{-2}$ by
$t \approx 1$, with the direction of the score preserved far better than
its magnitude.

Third, learning: the transition law is estimated from noised sequences
alone by expectation--maximisation, with Fisher's identity making each
E-step gradient exact through one forward--backward sweep. At an equal
number of clean trajectories --- and against networks that, unlike the
estimator, are additionally given those clean trajectories and redraw the
corruption at every gradient step --- it attains 6.6--10.9$\times$ lower
schedule-pooled relative score error than a denoising score-matching
network given no structural information, and 2.3--6.1$\times$ against a
baseline given locality and weight sharing, on a four-size protocol; the boundaries of the advantage
are mapped rather than asserted --- added mixture capacity buys nothing
under the implemented grid and budget, a rank-one violation of the Markov
assumption is partly absorbed while a long-range one defeats the estimator
at coupling one tenth against the convolutional baseline, one fifth
against both, and every factor is stated under an
explicitly defined aggregation whose alternatives are recomputed in an
appendix.

The setting is deliberately solvable --- one-dimensional states, linear
dynamics, known channel --- and the claims are bounded by it. What the
solvable case supplies is a reference point: exactly what a diffusion score
is, what a structural assumption is worth where it holds exactly, and what
it costs to learn one when it must be learned.
